% 118-11(118쪽) — 곡선 y=4x^3 과 x축·두 직선 x=a, x=b 로 둘러싸인 두 부분 S_1(a~0) S_2(0~b). 스캔 화질이 낮아 TikZ 로 다시 그림.
% 원문에 있는 것만: 곡선 · 두 색칠한 부분과 라벨 S_1 S_2 · 라벨 a, b, O, x, y, y=4x^3. 눈금 숫자는 원문에 없다.
% a, b 의 자리는 원문 그림과 같은 비율(|a| 가 b 의 절반쯤)로만 잡았고 값은 적지 않는다.
\begin{tikzpicture}[x=1.45cm, y=0.42cm, line width=0.5pt]
  % S_1 (a ≤ x ≤ 0)
  \fill[gray!18] plot[smooth, tension=0.6] coordinates {
    (-0.50,-0.50) (-0.40,-0.26) (-0.30,-0.11) (-0.20,-0.03) (-0.10,-0.004) (0,0) } -- (-0.50,0) -- cycle;
  % S_2 (0 ≤ x ≤ b)
  \fill[gray!18] plot[smooth, tension=0.6] coordinates {
    (0,0) (0.10,0.004) (0.20,0.03) (0.30,0.11) (0.40,0.26) (0.50,0.50) (0.60,0.86)
    (0.70,1.37) (0.80,2.05) (0.90,2.92) (1.00,4.00) } -- (1.00,0) -- cycle;
  \draw[->, >=stealth, line width=0.7pt] (-0.92,0) -- (1.33,0) node[below=1pt, font=\small] {$x$};
  \draw[->, >=stealth, line width=0.7pt] (0,-2.30) -- (0,5.50) node[left=1pt, font=\small] {$y$};
  % 직선 x=b
  \draw[line width=0.55pt] (1.00,0) -- (1.00,4.00);
  % 곡선
  \draw[line width=0.6pt] plot[smooth, tension=0.6] coordinates {
    (-0.78,-1.90) (-0.70,-1.37) (-0.60,-0.86) (-0.50,-0.50) (-0.40,-0.26) (-0.30,-0.11)
    (-0.20,-0.03) (-0.10,-0.004) (0,0) (0.10,0.004) (0.20,0.03) (0.30,0.11) (0.40,0.26)
    (0.50,0.50) (0.60,0.86) (0.70,1.37) (0.80,2.05) (0.90,2.92) (1.00,4.00) (1.05,4.63) };
  % S_1 안내 화살표와 라벨
  \draw[->, >=stealth, line width=0.35pt] (-0.62,-0.72) -- (-0.38,-0.17);
  \node[anchor=east, font=\small] at (-0.62,-0.82) {$S_1$};
  \node[font=\small] at (0.86,1.25) {$S_2$};
  % 라벨
  \node[above=1pt, font=\small] at (-0.50,0.10) {$a$};
  \node[below left=0pt and -2pt, font=\small] at (0,0) {$\pt{O}$};
  \node[below=2pt, font=\small] at (1.00,0) {$b$};
  \node[anchor=east, font=\small] at (1.02,5.25) {$y=4x^3$};
\end{tikzpicture}
